

%% 
%% This is a LaTeX document generated by automatic conversion of a
%% Mathematica notebook using Mathematica 4.0.
%% 
%% This document uses special macros that are defined in a LaTeX 
%% package file with a name of the form notebook*.sty.  Appropriate 
%% commands for loading the package have been included in this document.
%% The LaTeX package files can be found in a directory with path:
%% 
%% /usr/local/mathematica/SystemFiles/IncludeFiles/TeX/
%% 
%% The LaTeX package used in this document may require that your 
%% LaTeX implementation support fonts other than the standard Computer 
%% Modern fonts that are used by LaTeX.  See the Additional Information 
%% section under the TeXSave entry in the Mathematica Reference Guide 
%% for more information.
%% 
%% Implementation-specific instructions for installing TeX-related files
%% can be found at this URL.
%% 
%% http://www.wolfram.com/support/FrontEnds/Export/TeX
%% 
%% 

\documentclass{article}
\usepackage{notebook2e, latexsym}


\begin{document}

\Title{DESIGNING SIMPLE FILTERS}
\Subtitle{LECTURE 20 }
\Subsubtitle{ELE314: LINEAR SIGNALS AND SYSTEMS}
\Subsubtitle{Mariusz Jankowski \copyright{}

University of Southern Maine

mjkcc@usm.maine.edu}
\dispSFinmath{<<\Mvariable{SignalPlot`}}
\begin{CellGroup}
\Section*{1 Introduction}


We will consider one application of Fourier transforms. We will discuss a simple approach to designing {\itshape finite-impulse response} (FIR) digital
filters. The method is known as {\itshape Least Mean Squared Error} (LMS) design or simply the "{\itshape window}" method. The method requires two
steps:

(i) given a desired filter specification in the frequency domain, compute the impulse response by way of the IDTFT,

(ii) truncate the infinite impulse response to a desired length by means of an appropriate window function.

 It is common to design a prototype lowpass filter and subsequently use frequency shifting to obtain band-pass or high-pass filters.




\end{CellGroup}
\begin{CellGroup}
\Section*{2 IDTFT method}


The goal of digital filter design is to obtain an impulse response \inlineTFinmath{h[n]} (here we restrict ourselves to FIR filters only), given
a prototype filter specification in the frequency domain, \inlineTFinmath{H(\omega )}. We thus need to make use of the inverse discrete-time Fourier
transform.



Lets assume we want to design a lowpass filter with a cut-off frequency of \inlineTFinmath{{{\omega }_c}=\multsp 0.4\multsp \pi }. The prototype
ideal filter has the following form:



\inlineTFinmath{H(\omega )\multsp =\multsp \big\{\MathBegin{MathArray}[c]{cc}
	1,&|\omega |\multsp \leq \multsp \multsp {{\omega }_c} \\
	0,&\Mvariable{else} 
  \MathEnd{MathArray}}

for all frequencies within the fundamental interval, \inlineTFinmath{-\pi \leq \omega \leq \pi }. This is known as the {\itshape zero-phase} ideal
lowpass filter.



\mathGraphic{lecture20.tex_gr1.eps}
It is a simple matter to calculate the impulse response from the given filter specification.

\begin{CellGroup}
\dispSFinmath{\frac{1}{2\pi }\int _{-0.4\pi }^{0.4\pi }{{\ExponentialE }^{\ImaginaryI \multsp \omega \multsp n}}\DifferentialD \omega \multsp //\Mfunction{ExpToTrig}}
\dispSFoutmath{\frac{\sin [1.25664\multsp n]}{n\multsp \pi }}

\end{CellGroup}
\dispSFinmath{\MathBegin{MathArray}{l}
h[0]=0.4\multsp ;  \\
\noalign{\vspace{0.9375ex}}  \\
h[\Mvariable{n\_}]:=\frac{\sin [n\multsp 0.4\multsp \pi ]}{n\multsp \pi }
\MathEnd{MathArray}}
\begin{CellGroup}
\dispSFinmath{\Muserfunction{LinePlot}[\Mfunction{Table}[\{n,h[n]\},\{n,-10,25\}],\multsp \Mvariable{PlotRange}->\Mvariable{All}];}
\mathGraphic{lecture20.tex_gr2.eps}

\end{CellGroup}
\begin{CellGroup}
\Exercise{Problem 1}
\ExerciseText{Obtain the impulse response of the ideal highpass filter with cutoff frequency at \inlineTFinmath{{{\omega }_c}=\multsp 0.4\multsp
\pi }.}



\end{CellGroup}
The ideal impulse response is infinite and non-casual, thus cannot be used to realize a practical filter. Two operations need to be performed to
make the filter practical in real applications: (i) truncation of the filter's impulse response to a desired length \inlineTFinmath{L} and (ii) shift
of the filter by \inlineTFinmath{\frac{L-1}{2}} (for odd \inlineTFinmath{L}) samples to achieve causality. As a result of step (i) we get the following
finite length filter.



\begin{CellGroup}
\dispSFinmath{\MathBegin{MathArray}{l}
L=11;  \\
\noalign{\vspace{0.9375ex}}  \\
\Mvariable{fir}=\Mfunction{Table}\big[h[n],\big\{n,-\frac{L-1}{2},\frac{L-1}{2}\big\}\big]//\Mfunction{Chop}
\MathEnd{MathArray}}
\dispSFoutmath{\MathBegin{MathArray}{l}
\{0,-0.0756827,-0.062366,0.0935489,0.302731,  \\
\noalign{\vspace{0.5ex}}
\hspace{1.em} 0.4,0.302731,0.0935489,-0.062366,-0.0756827,0\}\\
\MathEnd{MathArray}}

\end{CellGroup}
We now verify the frequency response of the finite length, casual filter {\itshape fir }(note the implied shift in the calculation of the frequency
response).



\dispSFinmath{H[\Mvariable{\omega \_}]:=\sum _{i=0}^{L-1}\Muserfunction{fir}[[i+1]]\multsp \exp [-\imag \multsp \omega \multsp i]}
The magnitude response of the filter is shown here. We use a decibel scale for the magnitude and estimate the peak stopband ripple.

\begin{CellGroup}
\dispSFinmath{\Mfunction{Plot}[\{-20,20\multsp \log [10,\multsp \Mfunction{Abs}[H[\omega ]/H[0]]]\},\multsp \{\omega ,0,\pi \}];}
\mathGraphic{lecture20.tex_gr3.eps}

\end{CellGroup}
Here we compare the actual magnitude response with the response of the ideal lowpass filter.

\begin{CellGroup}
\dispSFinmath{\MathBegin{MathArray}{l}
\Mfunction{Plot}[\{\Mfunction{If}[\omega <0.4\pi ,1,0],\Mfunction{Abs}[H[\omega ]/H[0]]\},  \\
\noalign{\vspace{0.5ex}}
\hspace{1.em} \{\omega ,0,\pi \},\Mvariable{PlotStyle}->\{\Mfunction{Hue}[0],\Mfunction{Hue}[2/3]\}];\\
\MathEnd{MathArray}}
\mathGraphic{lecture20.tex_gr4.eps}

\end{CellGroup}
\begin{CellGroup}
\Exercise{Problem 2}
\ExerciseText{Obtain the frequency response (magnitude and phase) of the casual, lowpass filter of length \inlineTFinmath{3L}. Compare with the ideal
response. Estimate the peak stopband ripple.}



\end{CellGroup}

\end{CellGroup}
\begin{CellGroup}
\Section*{3 Improving Filter Characteristics by Modifying the Transition Region}


We now investigate a method of improving filter performance. Specifically, we want to decrease the filter's stopband ripple. One simple approach
is to introduce a smooth transition region into our prototype filter. Here we use a raised cosine.

	\inlineTFinmath{H(\omega )\multsp =\multsp \big\{\MathBegin{MathArray}[c]{ll}
	1&\omega <{{\omega }_p} \\
	\frac{1}{2}\big(1+\cos\big(\pi \multsp \frac{\omega -{{\omega }_p}}{{{\omega }_s}-{{\omega }_p}}\big)\big),&{{\omega }_p}\leq \omega \leq {{\omega
}_s} \\
	0,&\omega >{{\omega }_s} 
  \MathEnd{MathArray}},



where \inlineTFinmath{{{\omega }_p}} and \inlineTFinmath{{{\omega }_s}} are the passband and stopband edge frequencies, respectively.  We are given
the following edge frequencies:



\dispSFinmath{{{\omega }_p}=0.3\multsp \pi ;\multsp {{\omega }_s}=0.5\pi ;}
Here are the frequency response formulas for the raised cosine lowpass filter.

\dispSFinmath{\Mfunction{Clear}[H];}
\dispSFinmath{\MathBegin{MathArray}{l}
H[\Mvariable{\omega \_}]:=1\multsp /;\multsp \Mfunction{Abs}[\omega ]<0.3\multsp \pi ;  \\
\noalign{\vspace{0.9375ex}}  \\
H[\Mvariable{\omega \_}]:=\frac{1}{2}(1+\cos [5(\multsp \omega -0.3\multsp \pi )])\multsp /;(\multsp 0.3\multsp \pi \leq \omega \leq 0.5\pi \multsp
);  \\
\noalign{\vspace{1.17708ex}}  \\
H[\Mvariable{\omega \_}]:=\frac{1}{2}(1+\cos [5\multsp (\omega +0.3\pi )])\multsp /;(\multsp -0.5\multsp \pi \leq \omega \leq -0.3\pi \multsp );
 \\
\noalign{\vspace{0.90625ex}}  \\
H[\Mvariable{\omega \_}]:=0\multsp /;\multsp \Mfunction{Abs}[\omega ]>0.5\pi 
\MathEnd{MathArray}}
Here we plot the ideal lowpass filter and the "raised cosine" prototype.

\begin{CellGroup}
\dispSFinmath{\MathBegin{MathArray}{l}
\Mfunction{Plot}[\{\Mfunction{If}[\Mfunction{Abs}[\omega ]<0.4\pi ,1,0],H[\omega ]\},  \\
\noalign{\vspace{0.5ex}}
\hspace{1.em} \{\omega ,-\pi ,\pi \},\Mvariable{PlotStyle}->\{\Mfunction{Hue}[0],\Mfunction{Hue}[2/3]\}];\\
\MathEnd{MathArray}}
\mathGraphic{lecture20.tex_gr5.eps}

\end{CellGroup}
We now compute the IDTFT of the prototype filter. In this example we will use \buttonbox[RefGuideLink]{NIntegrate} to obtain the impulse response
of the desired filter. As in Problem 2, we desire a finite impulse response filter of length \inlineTFinmath{L=33}. As is clear from the above figure,
the frequency response is defined in terms of three sections. We will calculate the IDTFT (and immediately truncate) for each of the sections separately
and add the individual responses to obtain the desired finite length response. 



\begin{CellGroup}
\dispSFinmath{\MathBegin{MathArray}{l}
\Mfunction{Table}\big[\frac{1}{2\pi }\Big(\Mfunction{NIntegrate}[\exp [\imag \multsp \omega \multsp n],\{\omega ,-0.3\pi ,0.3\pi \}]+\Mvariable{NIntegrate}\big[
 \\
\noalign{\vspace{1.27083ex}}
\hspace{6.em} \frac{1}{2}(1+\cos [5(\multsp \omega +0.3\multsp \pi )])\multsp \exp [\imag \multsp \omega \multsp n],\{\omega ,-0.5\pi ,-0.3\pi \}\big]+
 \\
\noalign{\vspace{1.17708ex}}
\hspace{5.em} \Mfunction{NIntegrate}\big[\frac{1}{2}(1+\cos [5(\multsp \omega -0.3\multsp \pi )])\multsp \exp [\imag \multsp \omega \multsp n], 
\\
\noalign{\vspace{1.26042ex}}
\hspace{6.em} \{\omega ,0.3\pi ,0.5\pi \}\big]\Big),\{n,0,16\}\big]//\Mfunction{Chop}\\
\MathEnd{MathArray}}
\Message{\inlineSFinmath{\Mfunction{NIntegrate}::"ncvb"}\inlineSFinmath{:\multsp }NIntegrate failed to converge to prescribed accuracy after 7 recursive
bisections in \(\omega \) near \(\omega \) \(=\) \inlineSFinmath{-0.942478}}


\dispSFoutmath{\MathBegin{MathArray}{l}
\{0.4,0.29991,0.0900984,-0.0572778,-0.0649645,0,  \\
\noalign{\vspace{0.5ex}}
\hspace{1.em} 0.0354352,0.0163651,-0.0121286,-0.0142814,0,0.00681615,  \\
\noalign{\vspace{0.5ex}}
\hspace{1.em} 0.00264995,-0.00146866,-0.00097691,0,-0.000632771\}\\
\MathEnd{MathArray}}

\end{CellGroup}
\begin{CellGroup}
\dispSFinmath{d=\Mfunction{Join}[\Mfunction{Reverse}[\%],\Mfunction{Rest}[\%]]}
\dispSFoutmath{\MathBegin{MathArray}{l}
\{-0.000632771,0,-0.00097691,-0.00146866,0.00264995,0.00681615,  \\
\noalign{\vspace{0.5ex}}
\hspace{1.em} 0,-0.0142814,-0.0121286,0.0163651,0.0354352,0,-0.0649645,  \\
\noalign{\vspace{0.5ex}}
\hspace{1.em} -0.0572778,0.0900984,0.29991,0.4,0.29991,0.0900984,-0.0572778,  \\
\noalign{\vspace{0.5ex}}
\hspace{1.em} -0.0649645,0,0.0354352,0.0163651,-0.0121286,-0.0142814,0,  \\
\noalign{\vspace{0.5ex}}
\hspace{1.em} 0.00681615,0.00264995,-0.00146866,-0.00097691,0,-0.000632771\}\\
\MathEnd{MathArray}}

\end{CellGroup}
We now verify the frequency response of the "raised cosine", finite length, casual lowpass filter.

\dispSFinmath{\Mfunction{Clear}[H]}
\dispSFinmath{H[\Mvariable{\omega \_}]:=\sum _{i=0}^{32}d[[i+1]]\multsp \exp [-\imag \multsp \omega \multsp i]}
The magnitude response of the filter is shown here. We use a dB scale for the magnitude and estimate the peak stopband ripple.

\begin{CellGroup}
\dispSFinmath{\Mfunction{Plot}[\{-44,20\multsp \log [10,\multsp \Mfunction{Abs}[H[\omega ]/H[0]]]\},\multsp \{\omega ,0,\pi \}];}
\mathGraphic{lecture20.tex_gr6.eps}

\end{CellGroup}
\begin{CellGroup}
\Exercise{Problem 3}
\ExerciseText{Compare the frequency domain characteristics of the raised cosine filter with those of the filter obtained in Problem 2. Specifically
consider peak attenuation and width of transition region. Plot the two frequency responses on one diagram.}

\end{CellGroup}

\end{CellGroup}
\begin{CellGroup}
\Section*{4 Improving Filter Characteristics by the "Window" Method}


We now discuss a second method of improving the filter characteristics. The truncation of the impulse response of an ideal filter leads to the well
known Gibbs' phenomenon and poor stopband attenuation. Smooth window functions, such as Bartlett, von Hann, Hamming and Blackman, may be used to
improve filter performance (at the cost of a wider transition band). "Windowing" is an operation of multiplying a signal \inlineTFinmath{x[n]} by
a finite-duration {\itshape window signal} \inlineTFinmath{w[n]}. We get,



	\inlineTFinmath{s[n]\multsp =\multsp \multsp x[n]\multsp w[n]}



Simple truncation, as used in earlier examples, is equivalent to using the rectangular window. Here are the basic fixed window functions.

	Window name			Window function, \inlineTFinmath{w[n]}, \inlineTFinmath{-M\leq n\leq M}

	Rectangular			\inlineTFinmath{\frac{1}{2\multsp M+1}}

	Bartlett			\inlineTFinmath{1-\frac{|n|}{M+1}}

	von Hann			\inlineTFinmath{\frac{1}{2}\Big(1+\cos\Big(\frac{2\pi \multsp n}{2\multsp M+1}\Big)\Big)}

	Hamming			\inlineTFinmath{0.54+0.46\multsp \cos\Big(\frac{2\pi \multsp n}{2\multsp M+1}\Big)}

	Blackman			\inlineTFinmath{0.42+0.5\multsp \cos\Big(\frac{2\pi \multsp n}{2M+1}\Big)+0.08\multsp \cos\Big(\frac{4\pi \multsp n}{2M+1}\Big)}



Here is an example of the von Hann window.

\dispSFinmath{\Muserfunction{hann}[\Mvariable{n\_},\Mvariable{M\_}]:=\Mfunction{If}\big[\Mfunction{Abs}[n]\leq M,\frac{1}{2}\Big(1+\cos \big[\frac{2\pi
\multsp n}{2\multsp M+1}\big]\Big),0\big]}
\dispSFinmath{\Muserfunction{LinePlot}[\Mfunction{Table}[\{n,\Muserfunction{hann}[n,16]\},\{n,-20,20\}]];}
We now use the Hann window to window the impulse response of the example ideal lowpass filter.

\dispSFinmath{s[\Mvariable{n\_}]:=\multsp h[n]\multsp \Muserfunction{hann}[n,\multsp 16]}
\begin{CellGroup}
\dispSFinmath{\Mvariable{fir}=\Mfunction{Table}[s[n],\multsp \{n,-16,16\}]\multsp //\Mfunction{Chop}}
\dispSFoutmath{\MathBegin{MathArray}{l}
\{0.0000428371,0,-0.00120189,-0.00153958,0.00269062,0.00688024,  \\
\noalign{\vspace{0.5ex}}
\hspace{1.em} 0,-0.0144249,-0.01225,0.0165148,0.0357075,0,-0.0652284,  \\
\noalign{\vspace{0.5ex}}
\hspace{1.em} -0.0574158,0.0901984,0.299995,0.4,0.299995,0.0901984,-0.0574158,  \\
\noalign{\vspace{0.5ex}}
\hspace{1.em} -0.0652284,0,0.0357075,0.0165148,-0.01225,-0.0144249,0,  \\
\noalign{\vspace{0.5ex}}
\hspace{1.em} 0.00688024,0.00269062,-0.00153958,-0.00120189,0,0.0000428371\}\\
\MathEnd{MathArray}}

\end{CellGroup}
\dispSFinmath{S[\Mvariable{\omega \_}]:=\sum _{i=0}^{32}\Muserfunction{fir}[[i+1]]\multsp \exp [-\imag \multsp \omega \multsp i]}
The frequency response of the example filter designed using the length \inlineTFinmath{L=33} von Hann window is shown below.



\begin{CellGroup}
\dispSFinmath{\Mfunction{Plot}[\{-44,20\multsp \log [10,\multsp \Mfunction{Abs}[S[\omega ]/S[0]]]\},\multsp \{\omega ,0,\pi \}];}
\mathGraphic{lecture20.tex_gr7.eps}

\end{CellGroup}

\end{CellGroup}
\begin{CellGroup}
\Section*{5 Digital Filtering of Audio Data}


In this section, the filters designed in the earlier sections will be used to process voice data. The processing is implemented, as is typical, using
time-domain convolution. Here is a simple convolution code. This sets-up the path to the data directory.

\dispSFinmath{\MathBegin{MathArray}{l}
\Mvariable{dataDir}=\Mfunction{If}[\$System==\Mvariable{Linux},  \\
\noalign{\vspace{0.5ex}}
\hspace{2.em} /home/export/courses/ele314/data/,\backslash \backslash eelinux\backslash courses\backslash ele314\backslash data\backslash ];\\
\MathEnd{MathArray}}
Here we load an example signal.

\dispSFinmath{\Mvariable{sig}=\Muserfunction{Import}[\Mvariable{dataDir}<>\Mvariable{voice[1].wav}]\LeftDoubleBracket 1,1\RightDoubleBracket ;}
 The example signal had been sampled at a SampleRate of 11025 samples/sec.

\dispSFinmath{\Mfunction{SetOptions}[\Mvariable{ListPlay},\multsp \Mvariable{SampleRate}->11025,\multsp \Mvariable{PlayRange}->\{-1,1\}];}
Here we compare the original and filtered signals.

\begin{CellGroup}
\dispSFinmath{\Mfunction{ListPlay}[\Mfunction{Join}[\Mvariable{sig},\Mfunction{ListConvolve}[\Mvariable{fir},\Mvariable{sig}]]];}
\mathGraphic{lecture20.tex_gr8.eps}

\end{CellGroup}

\end{CellGroup}
\begin{CellGroup}
\Section*{PROBLEM SOLUTIONS}
\begin{CellGroup}
\Exercise{Problem 1}
\ExerciseText{Obtain the impulse response of the ideal highpass filter with cutoff frequency at \inlineTFinmath{{{\omega }_c}=\multsp 0.4\multsp
\pi }.}


{\bfseries Solution}. Using the IDTFT we get the expression for the impulse response of the ideal highpass filter.



\begin{CellGroup}
\dispSFinmath{\frac{1}{2\pi }\big(\int _{-\pi }^{-0.4\pi }{{\ExponentialE }^{\ImaginaryI \multsp \omega \multsp n}}\DifferentialD \omega \multsp
+\int _{0.4\pi }^{\pi }{{\ExponentialE }^{\ImaginaryI \multsp \omega \multsp n}}\DifferentialD \omega \big)//\Mfunction{ExpToTrig}}
\dispSFoutmath{-\frac{\sin [1.25664\multsp n]}{n\multsp \pi }+\frac{\sin [n\multsp \pi ]}{n\multsp \pi }}

\end{CellGroup}

\end{CellGroup}
\begin{CellGroup}
\Exercise{Problem 2}
\ExerciseText{Obtain the frequency response (magnitude and phase) of the length \inlineTFinmath{3L}, casual lowpass filter. Compare with the ideal
response. Estimate the peak stopband ripple.}


\ExerciseText{{\bfseries Solution}.}


\begin{CellGroup}
\dispSFinmath{\MathBegin{MathArray}{l}
L=33;  \\
\noalign{\vspace{0.9375ex}}  \\
\Mvariable{fir}=\Mfunction{Table}\big[h[n],\big\{n,-\frac{L-1}{2},\frac{L-1}{2}\big\}\big]//\Mfunction{Chop}
\MathEnd{MathArray}}
\dispSFoutmath{\MathBegin{MathArray}{l}
\{0.0189207,0,-0.0216236,-0.0143921,0.0155915,0.027521,0,  \\
\noalign{\vspace{0.5ex}}
\hspace{1.em} -0.0336367,-0.0233872,0.0267283,0.0504551,0,-0.0756827,  \\
\noalign{\vspace{0.5ex}}
\hspace{1.em} -0.062366,0.0935489,0.302731,0.4,0.302731,0.0935489,-0.062366,  \\
\noalign{\vspace{0.5ex}}
\hspace{1.em} -0.0756827,0,0.0504551,0.0267283,-0.0233872,-0.0336367,  \\
\noalign{\vspace{0.5ex}}
\hspace{1.em} 0,0.027521,0.0155915,-0.0143921,-0.0216236,0,0.0189207\}\\
\MathEnd{MathArray}}

\end{CellGroup}
\dispSFinmath{\Muserfunction{H33}[\Mvariable{\omega \_}]:=\Muserfunction{H33}[\omega ]=\sum _{i=0}^{L-1}\Muserfunction{fir}[[i+1]]\multsp \exp [-\imag
\multsp \omega \multsp i]}
The magnitude response of the filter is shown here. We use a dB scale for the magnitude and estimate the peak stopband ripple. The stopband attenuation
has not changed.

\begin{CellGroup}
\dispSFinmath{\Mfunction{Plot}[\{-20,20\multsp \log [10,\multsp \Mfunction{Abs}[\Muserfunction{H33}[\omega ]/\Muserfunction{H33}[0]]]\},\multsp \{\omega
,0,\pi \}]}
\mathGraphic{lecture20.tex_gr9.eps}
\dispSFoutmath{\SkeletonIndicator \Mvariable{Graphics}\SkeletonIndicator }

\end{CellGroup}

\end{CellGroup}
\begin{CellGroup}
\Exercise{Problem 3}
\ExerciseText{Compare the frequency domain characteristics of the raised cosine filter with those of the filter obtained in Problem 2. Specifically
consider peak attenuation and width of transition region. Plot the two frequency responses on one diagram.}
{\bfseries Solution}. Here is a plot of the frequency responses of the two filters.



\begin{CellGroup}
\dispSFinmath{\MathBegin{MathArray}{l}
\Mfunction{Plot}[\{20\multsp \log [10,\Mfunction{Abs}[H[\omega ]/H[0]]],20\multsp \log [10,\Mfunction{Abs}[\Muserfunction{H33}[\omega ]/\Muserfunction{H33}[0]]]\},\multsp
  \\
\noalign{\vspace{0.5ex}}
\hspace{1.em} \{\omega ,0,\pi \},\Mvariable{PlotStyle}->\{\Mfunction{Hue}[0],\Mfunction{Hue}[2/3]\}]\\
\MathEnd{MathArray}}
\mathGraphic{lecture20.tex_gr10.eps}
\dispSFoutmath{\SkeletonIndicator \Mvariable{Graphics}\SkeletonIndicator }

\end{CellGroup}

\end{CellGroup}

\end{CellGroup}
\begin{CellGroup}
\Section*{HOMEWORK PROBLEMS}
\begin{CellGroup}
\Subsubsection*{Homework Problem 20.1}
(a) Obtain the impulse response of a filter defined by the following equations:

\dispSFinmath{\MathBegin{MathArray}{l}
H[\Mvariable{\omega \_}]:=1\multsp /;\multsp \Mfunction{Abs}[\omega ]<0.3\multsp \pi ;  \\
\noalign{\vspace{0.5ex}}  \\
H[\Mvariable{\omega \_}]:=(0.5\multsp \pi -\omega )/(0.5\multsp \pi -0.3\multsp \pi \multsp )\multsp /;(\multsp 0.3\multsp \pi \leq \omega \leq 0.5\pi
\multsp );
\MathEnd{MathArray}}
\dispSFinmath{\MathBegin{MathArray}{l}
H[\Mvariable{\omega \_}]:=(\omega +0.5\multsp \pi )/(0.5\multsp \pi -0.3\multsp \pi \multsp )\multsp /;(\multsp -0.5\multsp \pi \leq \omega \leq
-0.3\pi \multsp );  \\
\noalign{\vspace{0.5ex}}  \\
H[\Mvariable{\omega \_}]:=0\multsp /;\multsp \Mfunction{Abs}[\omega ]>0.5\pi 
\MathEnd{MathArray}}
\begin{CellGroup}
\dispSFinmath{\Mfunction{Plot}[H[\omega ],\{\omega ,-\pi ,\pi \}];}
\mathGraphic{lecture20.tex_gr11.eps}
(b) Obtain the frequency response of the filter and compare with the ideal lowpass and "raised cosine" of Section 3. Specifically compare width of
transition region and peak stopband attenuation.


\end{CellGroup}

\end{CellGroup}
\begin{CellGroup}
\Subsubsection*{Homework Problem 20.2}
Compare the frequency response of the L\(=\)33 windows of type Bartlett, Hann, and Blackman. Which one has the narrowest transition region, which
one gives the best stopband attenuation. 


\end{CellGroup}
\begin{CellGroup}
\Subsubsection*{Homework Problem 20.3}
Design a half-band high-pass FIR filter with stopband attenuation \inlineTFinmath{{A_s}\leq \multsp 50\multsp \Mvariable{dB}} and transition width
 \inlineTFinmath{\Mvariable{\Delta \omega }\leq \multsp 0.12\multsp \pi }. Filter the example audio signal using the designed filter, comment on
the result.




\end{CellGroup}

\end{CellGroup}

\end{document}
